% Werke X/1 (1917), native234–236, printed227–229.
\begingroup
\clearpage\SourceLeaf{234}\ArithmeticHead{227}{R}{\scriptsize DARSTELLUNG DER ELLIPTISCHEN TRANSZENDENTEN.}
\NumberedItem{10}
\HistoricCitation{[Ein Zettel in Fi Nr. 6, Kapsel 50.]}
\par\indent Setzt man
\[\begin{aligned}
\operatorname{tg}T\surd\frac ab&=\operatorname{tg}T'\\[.5em]
\operatorname{tg}2T'\surd\frac{a'}{b'}&=\operatorname{tg}2T''\\[.5em]
\operatorname{tg}4T''\surd\frac{a''}{b''}&=\operatorname{tg}4T'''\\[.4em]
&\text{u.s.w.}
\end{aligned}\]
\noindent so wird
\[\begin{aligned}
\int\frac{dT}{\surd(aa\sin T^2+bb\cos T^2)}
&=\int\frac{dT'}{\surd(a'a'\sin2{T'}^2+b'b'\cos2{T'}^2)}\\[.7em]
&=\int\frac{dT''}{\surd(a''a''\sin2{T''}^2+b''b''\cos2{T''}^2)}\mathpunct{\raisebox{.5ex}{,}}
\end{aligned}\]
\noindent auch ist
\[\frac{\sin2T}{\surd(aa\sin T^2+bb\cos T^2)}
=\frac{\sin2T'}{\surd(a'a'\sin2{T'}^2+b'b'\cos2{T'}^2)}\mathpunct{\raisebox{.5ex}{.}}\]
\HistoricCitation{[Zwei Zettel in Fi Nr. 6, Kapsel 50.]}
\noindent [Erster Zettel]
\NumberedItem{11}
\[\begin{aligned}
a&=\mu(1+2xx+2x^8+2x^{18}+\text{etc.})^2=\mu fx\\[.5em]
b&=\mu(1-2xx+2x^8-2x^{18}+\text{etc.})^2=\mu fix\\[.5em]
c&=\mu(x^{\frac12}+x^{\frac92}+x^{\frac{25}2}+\text{etc.})^2=\mu\ldots
\end{aligned}\]
\[aa-bb=16cc\]
\[a+b=2a',\quad ab=b'b',\quad cc=a'c',\quad a-b=8c'\text{ etc.}\]
\noindent \(x\) ist die Grenze von
\[\frac ca,\quad\surd\frac{c'}{a'},\quad{}^4\!\surd\frac{c''}{a''},\quad{}^8\!\surd\frac{c'''}{a'''}\ldots\]
\noindent Ferner ist das ar.-geom.\ Mittel zwischen \(a\) und \(4c\)
\[=-\frac{\mu\pi}{2\log x}=\nu.\]
\par\bigskip\hfill\(29^*\)
\clearpage\SourceLeaf{235}\ArithmeticHead{228}{L}{ANALYSIS.\quad NACHLASS.}
\noindent Es ist also
\[\log x=-\tfrac12\frac{\mu\pi}{\nu}\mathpunct{\raisebox{.5ex}{.}}\qquad\operatorname{Log}\frac{k\pi}2=9{,}8339041884.\]
\NumberedItem{12}
\par\indent Setzt man
\[\surd(aa\sin T^2+bb\cos T^2)=\Delta,\]
\[\int\frac{dT}{\Delta}=\frac\theta\mu\mathpunct{\raisebox{.5ex}{,}}\]
\noindent so erhält man \(\theta\) durch folgenden Algorithmus
\[\begin{aligned}
\operatorname{tg}T\surd\frac ab&=\operatorname{tg}T'\\[.5em]
\operatorname{tg}2T'\surd\frac{a'}{b'}&=\operatorname{tg}2T''\\[.5em]
\operatorname{tg}4T''\surd\frac{a''}{b''}&=[\operatorname{tg}]4T'''\\[.5em]
\operatorname{tg}8T'''\surd\frac{a'''}{b'''}&=[\operatorname{tg}]8T^{\mathrm{IV}}\\[.4em]
&\text{etc.}
\end{aligned}\]
\noindent Dann ist \(\theta=T^\infty.\)
\par\indent Man hat ferner:
{\small
\[\operatorname{tg}\tfrac12T=
\frac{x^{\frac14}\sin\frac12\theta-x^{\frac94}\sin\frac32\theta+x^{\frac{25}4}\sin\frac52\theta-\text{etc.}}
{x^{\frac14}\cos\frac12\theta+x^{\frac94}\cos\frac32\theta+x^{\frac{25}4}\cos\frac52\theta+\text{etc.}}\]
\[\operatorname{tg}T'=
\frac{x^{\frac12}\sin\theta-x^{\frac92}\sin3\theta+x^{\frac{25}2}\sin5\theta-\text{etc.}}
{x^{\frac12}\cos\theta+x^{\frac92}\cos3\theta+x^{\frac{25}2}\cos5\theta+\text{etc.}}\]
\[\operatorname{tg}2T''=
\frac{x\sin2\theta-x^9\sin6\theta+x^{25}\sin10\theta-\text{etc.}}
{x\cos2\theta+x^9\cos6\theta+x^{25}\cos10\theta+\text{etc.}}\mathpunct{\raisebox{.5ex}{.}}\]}
\ClosingRule
\clearpage\SourceLeaf{236}\ArithmeticHead{229}{R}{ZUR RECTIFICATION DER ELLIPSE.}
\par\indent Eine independente Rechnung ist folgende
\[\frac ab\operatorname{tg}T=\operatorname{tg}U\]
{\small
\[\begin{array}{ll}
\operatorname{tg}T'=\surd\operatorname{tg}T.\operatorname{tg}U,&U'=\tfrac12(T+U)\\[.5em]
\operatorname{tg}2T''=\surd\operatorname{tg}2T'.\operatorname{tg}2U',&U''=\tfrac12(T'+U')\\[.5em]
\operatorname{tg}4T'''=\surd\operatorname{tg}4T''.\operatorname{tg}4U'',&U'''=\tfrac12(T''+U'')
\end{array}\]}
\begin{center}u.s.w.\end{center}
\noindent Dann ist
\[T^\infty=U^\infty=\theta\]
{\small
\[\begin{aligned}
&\int\surd(aa\sin T^2+bb\cos T^2)dT\\
&\quad=\frac\theta\mu\bigl(a'a'-2(a''a''-b''b'')-(4a'''a'''-b'''b''')\ldots\bigr)\\[.5em]
&\qquad+\tfrac12(a-b)\sin2U'-\tfrac12(a'-b')\sin4U''\\[.4em]
&\qquad-\tfrac12(a''-b'')\sin8U'''-\tfrac12(a'''-b''')\sin16U^{\mathrm{IV}}\ldots
\end{aligned}\]}
\noindent [Zweiter Zettel]
\HistoricTitle{Zur Rectification der Ellipse.}
\NumberedItem{13}
\par\indent Es sei
\[\int\surd(aa\cos T^2+bb\sin T^2)dT=W.\]
\par\indent Man bilde die Reihe der arithm.-geom.\ Mittel
\[\begin{array}{r|ll}
a,b&aa-bb=16cc&cc=a'c'\\[.4em]
a',b'&a'a'-b'b'=16c'c'&c'c'=a''c''\\[.4em]
a'',b''&a''a''-b''b''=16c''c''&c''c''=a'''c'''\\[.4em]
\text{etc.}&\text{etc.}&\text{etc.}\\[-.4em]
\begin{array}{c}.\\[-.8em].\\[-.8em].\end{array}&&\\
{}[M(a,b)=]\quad\alpha&&
\end{array}\]
\noindent Man hat dann ferner zu setzen
\[\begin{array}{lll}
\dfrac{4\sin T}a=p&cp=\sin2u&p'=\dfrac p{\cos u^2}\\[.8em]
\dfrac{4\sin T'}{a'}=p'&c'p'=\sin2u'&p''=\dfrac{p'}{\cos{u'}^2}\\[.8em]
\dfrac{4\sin T''}{a''}=p''&c''p''=\sin2u''&\text{etc.}\\
\begin{array}{c}.\\[-.8em].\\[-.8em].\\[-.8em].\end{array}&\begin{array}{c}.\\[-.8em].\\[-.8em].\\[-.8em].\end{array}&
\end{array}\]
\clearpage\endgroup
